\chapter{Introduction}
\label{ch:introduction}

\section{Background and Motivation}

Institutional policies govern the behavior of organizations and their members. Universities, corporations, and government agencies maintain extensive policy documents that specify obligations, permissions, and prohibitions for various stakeholders. However, these policies are typically written in natural language, making automated compliance verification challenging.

The Asian Institute of Technology (AIT) maintains numerous Policies and Procedures (P\&P) documents covering student conduct, financial obligations, accommodation rules, and grievance procedures. Manually interpreting and enforcing these policies is time-consuming and error-prone.

Recent advances in Large Language Models (LLMs) offer promising approaches for natural language understanding tasks, including text classification and information extraction. Meanwhile, formal logic provides a rigorous foundation for representing and reasoning about normative statements.

\section{Problem Statement}

The key challenges addressed in this thesis are:

\begin{enumerate}
    \item How can policy rules be automatically identified and extracted from natural language documents?
    \item What formal representation is appropriate for capturing the semantics of policy rules?
    \item How can formalized policies be translated into executable validation constraints?
\end{enumerate}

\section{Research Questions}

\begin{itemize}
    \item \textbf{RQ1}: Can Large Language Models effectively identify and classify policy rules in academic documents?
    \item \textbf{RQ2}: Is First-Order Logic with deontic extensions sufficient for formalizing institutional policies, or are higher-order logics required?
    \item \textbf{RQ3}: Can formalized policy rules be systematically translated to SHACL shapes for RDF-based validation?
\end{itemize}

\section{Objectives}

The objectives of this research are:

\begin{enumerate}
    \item Evaluate multiple LLM models for policy rule identification
    \item Develop a formalization approach using First-Order Logic with deontic operators
    \item Create a translation mechanism from FOL to SHACL shapes
    \item Validate the approach using AIT policy documents as a case study
\end{enumerate}

\section{Scope and Limitations}

This research focuses on:
\begin{itemize}
    \item English-language policy documents from AIT
    \item Declarative policy rules (obligations, permissions, prohibitions)
    \item Static analysis without runtime enforcement
\end{itemize}

Excluded from scope:
\begin{itemize}
    \item Multilingual policy documents
    \item Procedural or workflow policies
    \item Real-time compliance monitoring
\end{itemize}

\section{Contributions}

The main contributions of this thesis are:

\begin{enumerate}
    \item A systematic evaluation of five LLM models for policy rule classification
    \item Evidence that FOL with deontic operators is sufficient for institutional policy formalization
    \item An automated pipeline from natural language policies to SHACL validation constraints
    \item A gold standard dataset of 97 annotated policy rules
\end{enumerate}

\section{Thesis Structure}

This thesis is organized as follows:

\begin{itemize}
    \item \textbf{Chapter 2}: Literature Review --- Background on formal logic, LLMs, and SHACL
    \item \textbf{Chapter 3}: Methodology --- Research design and evaluation approach
    \item \textbf{Chapter 4}: Implementation --- System architecture and components
    \item \textbf{Chapter 5}: Results --- Experimental findings and analysis
    \item \textbf{Chapter 6}: Discussion --- Implications and limitations
    \item \textbf{Chapter 7}: Conclusion --- Summary and future work
\end{itemize}
